\section{Dynamic Languages}

Dynamic languages can be generalised as languages that differ several beheviours of the language to runtime, that would otherwise occur during compile time for static languages.\footnote{\url{https://developer.mozilla.org/en-US/docs/Glossary/Dynamic_programming_language}}

For the purpose of this project, we will focus on dynamic languages in the context of Javascript. In particular, Javascript exhibits the following behaviours which contribute to its dynamic nature.

\subsection{Eval}

In Javascript, the \textit{eval} function takes a string as parameter and evaluates it as if it were Javascript source code, returning the result of evaluating the string.\footnote{\url{https://developer.mozilla.org/en-US/docs/Web/JavaScript/Reference/Global_Objects/eval}}


\lstinputlisting{./code_examples/ch01/eval.js}

In the example above, the argument string passed to \textit{eval} is interpreted as valid Javascript code and evaluates to 2.

This mechanism allows code to be changed at runtime, making Javascript a dynamic language. However, there has been considerable criticism against the usage of \textit{eval}. Most notably, \textit{eval} is unsafe\footnote{\url{https://developers.google.com/web/fundamentals/security/csp#eval_too}} and prevents an optimising compiler from conducting data flow optimisations\footnote{\url{https://developer.mozilla.org/en-US/docs/Web/JavaScript/Reference/Global_Objects/eval#never_use_eval!}}.

\subsection{Object Modification}

Objects in Javascript can be modified at runtime by adding or removing property fields and behaviour. Contrary to a more static language like Java where object instances are created based on the attributes and behaviour of the instantiating class, Javascript objects can add properties via the \textit{Object.defineProperty()}\footnote{\url{https://developer.mozilla.org/en-US/docs/Web/JavaScript/Reference/Global_Objects/Object/defineProperty}} and delete properties via the \textit{delete}\footnote{\url{https://developer.mozilla.org/en-US/docs/Web/JavaScript/Reference/Operators/delete}} operators.

Functions can be defined on objects at runtime in the following manner:

\lstinputlisting{./code_examples/ch01/defineFunction.js}

In the above example, \textit{fooFunc} is only attached to \textit{obj} at runtime. It would not be possible to conduct static analysis on the location of \textit{fooFunc} using the information available at compile time.

\subsection{Dynamic Types}

Types in Javascript are determined at runtime and can be altered at runtime. An AST\footnote{Abstract Syntax Tree} generated at compile time does not capture any information about the type of the variable.

In the following example, an AST is generated for a variable declaration. It can be observed that no information is available of the type of variable \textit{a}. Its type can be determined using the \textit{typeof} operator at runtime only.

\lstinputlisting{./code_examples/ch01/varDecl.js}

\begin{center}
    \includegraphics[height=150pt]{./images/var_decl.eps}
\end{center}

In addition, the type of \textit{a} can be altered at runtime. It is possible to assign a string or function to \textit{a}, changing its type.

\lstinputlisting{./code_examples/ch01/varAssignment.js}

Thus, it would be impossible to conduct static analysis by relying on the type system of a dynamic language.
